\chapter{Introduzione}
\label{ch:introduzione}

I tool di sicurezza che operano in user-mode identificano operazioni sospette
ispezionando la \emph{call stack} al momento in cui avvengono: creazione di
processi, allocazioni di memoria eseguibile, apertura di handle privilegiati.
Questa ispezione risale la catena di frame pointer e return address per
ricostruire la sequenza di chiamate che ha portato all'operazione. Se la call
chain risultante corrisponde a funzioni di sistema legittime, l'evento è
classificato come benigno; in caso contrario, viene segnalato come anomalo.

Il \emph{call stack spoofing} è una tecnica di \emph{defense evasion} che
sfrutta questa dipendenza: falsificando i frame pointer e i return address
visibili agli stack walker user-mode, è possibile presentare una call chain
apparentemente legittima anche durante operazioni sospette. L'efficacia della
tecnica dipende dall'architettura del sistema e dai meccanismi con cui il
sistema operativo gestisce la catena di chiamate.

Un caso emblematico è il \emph{worm}: una volta ottenuta l'esecuzione su un
sistema, deve allocare memoria eseguibile, creare processi per replicarsi e
comunicare con l'esterno — operazioni intrisecamente sospette che i tool di
sicurezza in user-mode ispezionano sistematicamente. Il call stack spoofing
consente al worm di mascherare queste operazioni presentando una call chain
apparentemente legittima, riducendo la probabilità di rilevamento senza
modificare il comportamento funzionale.

\paragraph{Struttura del documento.}
Il Capitolo~\ref{ch:intro} introduce i fondamenti dell'architettura ARM64
su Windows necessari per comprendere le tecniche analizzate: calling convention,
catena \texttt{x29}/\texttt{x30}, segmento \texttt{.pdata} e Thread Environment
Block. Il Capitolo~\ref{ch:scenari} documenta cinque scenari di call stack
spoofing con crescente complessità, ciascuno accompagnato dalle evidenze
raccolte con WinDbg, System Informer e Sysmon. Il Capitolo~\ref{ch:difensiva}
analizza i limiti degli stack walker user-mode e le firme residue rilevabili.
Il Capitolo~\ref{ch:conclusioni} sintetizza il tradeoff tra complessità
offensiva e copertura difensiva.
